\subsection{Solving Quadratic Equations by the Square Root Property} 
If a quadratic equation is given in the appropriate form, we can also solve it using the \makeblank[1in] root property. (We'll say ``square root property'' here because we're interested specifically in index \makeblanksmall.)
\begin{Example}
Solve using the square root property.
\begin{lpic}{}{6}
\regtextleft{0.25}{-1}{$(x-4)^2=9$};
\regtextleft{0.25}{-5}{Check:};
\end{lpic}
\end{Example}
\noi We can solve quadratic equations using the square root property that could not be solved by factoring.
\begin{Example}
The equation $(x+1)^2=6$ cannot be solved by factoring.
\begin{lpic}{}{4}
\regtextleft{0.25}{-1}{$(x+1)^2=6$};
\end{lpic}
However, it can be solved using the square root property.
\begin{lpic}{}{7}
\regtextleft{0.25}{-1}{$(x+1)^2=6$};
\regtextleft{0.25}{-4}{Check:};
\end{lpic}
\end{Example}
\noi In order to use the square root property to solve, we have to write our equation in the form of a \makeblank[1in] equation rather than a \makeblank[1.5in] equation. Is this always possible?